% TOP_PROBLEM: {"problemId": "problem.polynomial-time-solvability-of-condon-s-simple-stochastic-games", "canonicalTitle": "Polynomial-Time Solvability of Condon's Simple Stochastic Games", "releaseRank": 478, "primaryDomain": "theoretical_computer_science", "primaryDomainLabel": "Theoretical computer science", "displayStatus": "open", "releaseStatus": "open"}
% !TeX program = lualatex
% Definition and sources only; this file is not an LLM solution attempt.
\documentclass[11pt]{article}
\usepackage[margin=25mm]{geometry}
\usepackage{fontspec}
\setmainfont{Latin Modern Roman}
\usepackage{amsmath,amssymb,mathtools}
\usepackage{unicode-math}
\setmathfont{Latin Modern Math}
\usepackage{xurl,hyperref}
\hypersetup{colorlinks=true,urlcolor=blue}
\setcounter{secnumdepth}{0}
\setlength{\parindent}{0pt}
\setlength{\parskip}{5pt}
\setlength{\emergencystretch}{3em}
\begin{document}
\section*{\texorpdfstring{Polynomial-Time Solvability of Condon's Simple Stochastic Games}{Polynomial-Time Solvability of Condon's Simple Stochastic Games}}\label{478}
\subsection{Definitions and mathematical statement}
A simple stochastic game is a finite directed graph whose nonterminal vertices belong to a maximizing player, a minimizing player, or chance. Players choose the outgoing edge at their own vertices; chance vertices use the given rational probabilities. Success means reaching a designated target, with nonreaching plays counted as failures. For initial vertex $s$, the value is
\[
 v(s)=\sup_\sigma\inf_\tau
 {\mathbb{P}}_s^{\sigma,\tau}(\text{eventually reach the target}).
\]
The question asks for a deterministic polynomial-time method to compute these values exactly, equivalently to decide whether $v(s)\ge q$ for a rational threshold $q$, in the standard finite rational-input convention. Binary decisions and transition encodings are part of the instance. An approximate answer with unspecified dependence on the precision or an algorithm exponential in input length does not settle the selected complexity question.
\par\smallskip\textit{Formulation and concept references:} \href{https://dblp.org/rec/journals/iandc/Condon92.html}{[S1]}.

\subsection{Short English statement}
Can the exact optimal winning probabilities of finite simple stochastic reachability games be computed in polynomial time?
\subsection{Sources}
\begin{itemize}
\item[S1]Formal statement, dblp record header for journals/iandc/Condon92.
\url{https://dblp.org/rec/journals/iandc/Condon92.html}.
\textit{Formulation source.}
\end{itemize}
\end{document}
